% page 203 of the facsimile (image p-202 of the scan)
% block 167.6 x 254.9 mm ; left margin 34.4 mm
\pgc{12.700mm}{0.000mm}{
\l{}
\l{}
\l{}
\l{}
\l{\cel{54}195}
\l{}
\l{\sou{glacizo}. (\sou{milito}). Pento mikr-angula qua, en fortifikajo, abutas}
\l{\cel{3}a la ruro cirkondanta. - DEFRS.}
\l{}
\l{\sou{gladiatoro}.\cel{2}(en\cel{1}\sou{la Roma antiqua}). Viro quan onu igis kombatar,}
\l{\cel{3}en cirko, kontre altra viri o kontre bestii feroca, por}
\l{\cel{3}amuzar la plebeyo spektanta. - DEFIRS.}
\l{}
\l{\sou{gladiolo.}\cel{2}(\sou{bot.}) Planto analoga ad irido. - L.gladiolus. DEFIL.}
\l{}
\l{\sou{glano}. I. (\sou{bot.}) Frukto di la querko, di qua la perikarpo, union-}
\l{\cel{3}ita a la semino, kovresas da involukro squamoza. - II. (\sou{anat}.)}
\l{\cel{3}Extremajo di la peniso, di la klitorido. - III. Tufo de silko,}
\l{\cel{3}lano, e c. uzata kom ornivo en la pasmenti. - DEFIS.}
\l{}
\l{\sou{glando}. (\sou{anat.}) Organo ek mikra utrikuli, qua sekrecas ula liquidi}
\l{\cel{3}aden la korpo animalala. - DEFIS.}
\l{}
\l{\sou{glaso}. La vazo ordinara, pedoza, (generale ek vitro), en qua onu}
\l{\cel{3}drinkas vino od aquo. - DE.}
\l{}
\l{\sou{glata}. Qua ne sentigas asperaji kande onu igas lu glitar sur sua}
\l{\cel{3}pelo. - DR.}
\l{}
\l{\sou{glaukomo}. (\sou{patol}.)Okulo-morbo, pro qua la humoro vitratra divenas}
\l{\cel{3}opaka, e la okul-fundo verdatra. - DEF.}
\l{}
\l{\sou{glaukonito}. (\sou{mineral.).}\cel{2}Hidroza silikato de fero, de alumino e}
\l{\cel{3}de kalio, quaza kreto verdatra. - DEFIS.}
\l{}
\l{\sou{glavo}. Espado tranchiva, kun lamo kurta e rekta, di qua la extre-}
\l{\cel{3}majo esas du-bizela. - FL.}
\l{}
\l{\sou{glebo}.\cel{2}Mikra maso de tero kompakta. - FI.}
\l{}
\l{gleno. (\sou{anat.}) Kavajo poke profunda di osto en qua esas inkastri-}
\l{\cel{3}ta altra osto.}
\l{}
\l{glezar. (\sou{trans.}) Aplikar (sur terakotajo, ceramikajo) induturo}
\l{\cel{3}vitrigebla. - DE.}
\l{}
\l{glicerio. (\sou{bot.)}\cel{2}Genero de planti perena, ek la familio \textquotedbl{}grami-}
\l{\cel{3}nei\textquotedbl{}. -L.\cel{1}\sou{glyceria}.}
\l{}
\l{glicerino.\cel{2}(\sou{kemio}).Korpo neutra, siropatra, quan onu obtenas per}
\l{\cel{3}la duimigo di grasajo (stearino, oleino, e c.) - DEFIRS.}
\l{}
\l{glicino. (\sou{bot.})\cel{2}Arbusto klimera, de la familio \textquotedbl{}papilionacei\textquotedbl{}.}
\l{\cel{3}- L.\cel{1}\sou{glycine.}\cel{2}- DEFIS.}
\l{}
\l{glikogena. (fiziol.)\cel{2}Qua produktas sukro.}
\l{}
\l{\cel{1}glikogenezo.\cel{2}(\sou{kemio biologiala)}.\cel{2}Formaceso di la sukro (ex.:}
\l{\cel{4}en la hepato).}
}
